\scanpage{8}
\mspage{7}
On dit que $F,F'$ sont \uline{compatibles} si les conditions en question sont \unsure{satisfaites}.

Pour ceci, il f et \unsure{s.|s} que pour \unsure{toute|tous} $F_0\in\Sigma'_F$, $F'_0\in\Sigma'_{F'}$, $F_0$ et $F''_0$ soient compatibles.

Donc, dans l'ens \inserted{ordonné $\mathcal{F}_{\text{él}}$} des figures élémentaires, on introduit une relation (symetrique, reflexive)
\[
\begin{aligned}
&\mathcal{R}\subset\mathcal{F}_{\text{él}}\times\mathcal{F}_{\text{él}}&&\deleted{\{\unsure{F,F'|F_0,F'_0}\}}\quad \text{$F_0$ comp. avec $F'_0$}\\
\text{ou encore}\quad&\mathcal{S}\in\mathcal{P}_2(\mathcal{F}_{\text{él}})&&\deleted{\text{\unsure{ou|cou}}}\quad\text{\uline{paires compatibles}}
\end{aligned}
\]
\uline{NB} Si \deleted{\unsure{$X$|$2$}} $F_0\unsure{\to|,} F'_0$ sont des fig. él. telles que $F_0\le F'_0$, alors elles sont compatibles. Mieux, si $F'_0,F''_0\le F$, alors $F_0$ et $F'_0$ sont compatibles entre elles

\uline{Configuration} Ens de figures deux : deux compatibles. Configuration élémentaire : ens des figures élémentaires deux 2 deux compatibles.
